\chapter{Continuous integration and quality}
\label{chap:cicd}
% Corresponds to phase 02 of the roadmap.

\section{Jenkins pipeline}
A Jenkins Declarative Pipeline builds and redeploys the dashboard automatically. It runs four stages in sequence -- checkout, test, build \& deploy, health check -- and, critically, the deployment stage is unreachable if the test stage fails: a broken change never reaches production. Triggering is handled by a GitHub webhook for near-instant deployment, with a 20-minute SCM poll kept as a fallback in case the webhook delivery ever fails silently.

The health check was adjusted once the dashboard port stopped being published on every interface (chapter~\ref{chap:securisation}): instead of calling the port from the host, the stage runs the request inside the container, up to six times five seconds apart, so that it stays valid whatever address the port is published on and tolerates the few seconds the application needs to start.

\section{Source control and credentials}
The dashboard's code lives in a private GitHub repository. Two different SSH deploy keys are used on purpose. The virtual machine holds a key with write access, so that a quick fix can be committed from the server itself. Jenkins holds its own key, restricted to read access, because it only needs to fetch the code: a compromised pipeline can therefore never modify the source. Secrets (dashboard credentials, Proxmox token, Discord webhook, Jenkins API token) are kept in a configuration file that Git ignores, so a deployment never overwrites them and they never appear in the repository.

\section{Automated tests}
The test suite, 24 tests at the time of writing, runs on Node.js's built-in test runner (\texttt{node --test}), which requires no extra dependency. It covers pure functions (SSH log line parsing, private-IP detection), the authentication middleware (including the anti-brute-force lockout behaviour described in chapter~\ref{chap:securisation}), and the API routes through short-lived HTTP servers started on ephemeral ports for the duration of a test. The Jenkins pipeline runs this suite inside an ephemeral, disposable \texttt{node:20-alpine} container rather than on the host directly, so the test environment matches a clean machine rather than whatever happens to already be installed on the VM.

\section{Verifying that backups actually work}
A backup that has never been restored is only a hypothesis. A dedicated script, run a few minutes after the nightly backup job, opens the freshly produced archive, checks that it is not corrupted, and confirms the application directory it is supposed to contain is actually present inside it -- then writes the result to a file the dashboard reads and displays, and sends a Discord alert if the check fails. The check runs fifteen minutes after the backup (02:00 then 02:15) and ends in one of three outcomes: valid, valid but incomplete (the project is present but secondary items such as firewall or SSH configuration are missing), or invalid (corrupted archive or missing project, meaning it could not be used to restore the service). One operational lesson came with it: since the result file is mounted into the dashboard container, it must exist as a real file before the first start, otherwise Docker creates an empty folder in its place and the mount breaks; the script is therefore run once by hand before the first deployment.

\section{Difficulties encountered}
Getting Jenkins to build from the exact same working directory as the intern's own manual deployments (needed so both share the same \texttt{data/} folder and container names) surfaced a chain of independent permission issues that had to be diagnosed one layer at a time: a DNS resolution fault on the VM unrelated to Jenkins itself (traced to \texttt{systemd-resolved} instability and fixed by forcing reliable DNS servers); Jenkins' own user lacking write access not just to the project folder but to its \emph{parent} directory, where Jenkins creates a sibling control directory; and finally Git's "dubious ownership" protection (CVE-2022-24765) refusing to operate in a directory owned by a different user than the one running the command, resolved by explicitly marking the directory as safe for Jenkins' user. Each layer produced a distinct, initially confusing error message, and the only way through was to reproduce the exact failing command directly as the \texttt{jenkins} user via \texttt{sudo -u jenkins}, rather than iterate through the Jenkins web UI.

\clearpage
